\documentclass{article}

\usepackage{neurips_2025}
\usepackage[utf8]{inputenc}
\usepackage[T1]{fontenc}
\usepackage{hyperref}
\usepackage{url}
\usepackage{booktabs}
\usepackage{amsfonts}
\usepackage{nicefrac}
\usepackage{microtype}
\usepackage{xcolor}
\usepackage{amsmath}
\usepackage{amssymb}
\usepackage{graphicx}
\usepackage{multirow}
\usepackage{enumitem}

\title{Matched-Budget Evaluation of Weak-View Residuals for One-Run Differential Privacy Auditing}

\author{Anonymous Authors}

\newcommand{\eps}{\varepsilon}
\newcommand{\del}{\delta}
\newcommand{\rclean}{r_{\mathrm{clean}}}
\newcommand{\rmean}{r_{\mathrm{mean}}}
\newcommand{\rworst}{r_{\mathrm{worst}}}
\newcommand{\ragree}{r_{\mathrm{agree}}}
\newcommand{\lambdasection}{\texorpdfstring{\(\lambda\)}{lambda}}

\begin{document}

\maketitle

\begin{abstract}
One-run privacy auditing makes empirical auditing of differentially private (DP) training practical, and recent work showed that quantile-regression calibration substantially strengthens final-model-only black-box auditing. We study a narrower follow-up question: after adopting that calibrated single-view baseline, do repeated weak, label-preserving views add measurable privacy evidence once every extra query is charged honestly? We evaluate this in a one-run audit with fixed audit pools, randomized orderings, and matched-total-query comparisons on Fashion-MNIST and CIFAR-10 under realized privacy budgets near \(\eps=3.94\) and \(\eps=7.95\). The weak-view method combines clean, worst-view, and agreement residuals derived from \(K=4\) weak views, while all residual calibrators and score-selection steps are fit only on public data. Across four seeds, 240 audit pairs per checkpoint, and query budgets from 40 to 1200, the structured weak-view score does not improve the primary metric, the empirical \(\eps\) lower bound at \(\del=10^{-5}\), over the clean quantile baseline in any setting. On CIFAR-10, for example, the clean quantile baseline reaches mean full-pool lower bounds of 0.136 at target \(\eps=4\) and 0.158 at target \(\eps=8\), versus 0.093 and 0.078 for the structured weak-view score. The main finding is therefore negative but operationally useful: in this compute-bounded one-run setting, extra weak-view queries do not yield additional privacy evidence beyond a strong single-view quantile baseline when evaluated under matched query accounting.
\end{abstract}

\section{Introduction}
Empirical privacy auditing asks whether a training pipeline behaves as privately as its formal accounting suggests. One-run auditing made this practical by embedding many neighboring comparisons into a single training run \citep{steinke2023privacy}, and subsequent work strengthened the conversion from empirical distinguishability to privacy evidence \citep{mahloujifar2024auditing}. In the final-model-only black-box setting, \citet{liu2025enhancing} further showed that the per-example score itself matters substantially: quantile-regression calibration can outperform raw loss and materially strengthen one-run auditing.

This paper studies a deliberately narrower question than prior multi-query membership inference work. Repeated queries, augmentations, and robustness-based scoring are already known to help membership inference in other settings \citep{he2022semileak,rahimian2020sampling,choquettechoo2021labelonly,wu2024yoqo}. What remains unclear is whether such extra views still help once the comparison point is not a weak single-query baseline, but rather the quantile-calibrated single-view one-run auditor of \citet{liu2025enhancing}, and once every additional query is charged explicitly.

We therefore evaluate weak-view residual auditing under a matched-budget protocol. The protocol fixes a finite audit pool, counts the clean query together with all weak-view queries, and compares methods as a function of consumed queries rather than candidate count. The resulting claim is modest: this is not a new attack class, but a protocol-focused empirical test of whether weak-view residuals add privacy evidence beyond a strong single-view baseline.

Our contributions are:
\begin{itemize}[leftmargin=*,topsep=2pt,itemsep=2pt]
    \item We instantiate a public-data-only weak-view residual auditor for one-run DP auditing, combining clean-view, worst-view, and agreement residuals without changing the underlying one-run privacy conversion.
    \item We evaluate that auditor under a matched-total-query protocol on Fashion-MNIST and CIFAR-10 with four seeds, 240 audit pairs per checkpoint, and query budgets from 40 to 1200.
    \item We find a consistent negative result: the structured weak-view score does not improve the empirical \(\eps\) lower bound over the clean quantile baseline in any of the four dataset/privacy settings, and proxy \(\lambda\)-selection frequently collapses to \(\lambda=0\).
\end{itemize}

Section~\ref{sec:related} reviews the most relevant auditing and repeated-query literature. Section~\ref{sec:method} describes the weak-view residual score and matched-budget protocol. Section~\ref{sec:experiments} presents the experimental setup and results. Section~\ref{sec:discussion} discusses implications and limitations, and Section~\ref{sec:conclusion} concludes.

\section{Related Work}
\label{sec:related}

\paragraph{One-run auditing and privacy conversion.}
Privacy auditing with one training run was introduced by \citet{steinke2023privacy}, who showed how to pack many neighboring comparisons into a single audited training job. \citet{mahloujifar2024auditing} strengthened the privacy conversion through \(f\)-DP analysis. Our work leaves the neighboring-dataset construction and empirical lower-bound conversion unchanged; we vary only the final-model score and the fairness protocol used to compare scores.

\paragraph{Quantile-regression one-run auditing.}
The closest baseline is \citet{liu2025enhancing}, which argues that final-model-only one-run auditing benefits substantially from quantile-regression calibration of black-box scores. Their result sharply raises the bar for follow-up work: extra signal sources must beat a strong calibrated baseline, not merely raw loss. Our paper is best understood as an incremental evaluation around that baseline.

\paragraph{Repeated-query and multi-view membership inference.}
Prior work already established that repeated queries can help membership inference in broader settings. Semi-Leak uses repeated augmented queries against semi-supervised models \citep{he2022semileak}. Sampling Attacks amplify label-only membership inference through repeated stochastic queries \citep{rahimian2020sampling}. Label-only attacks estimate robustness through many perturbation queries \citep{choquettechoo2021labelonly}, while YOQO reduces the query burden of that line of work \citep{wu2024yoqo}. These papers constrain our novelty claim, but they do not answer our question directly because they do not study one-run DP auditing with empirical \(\eps\) lower bounds, quantile-calibrated single-view baselines, and matched total-query accounting.

\paragraph{Other recent auditing directions.}
Recent auditing work tightened black-box auditing \citep{annamalai2024nearly}, introduced adversarial samples in final-model-only settings \citep{yoon2024adversarial}, avoided retraining through PANORAMIA \citep{kazmi2024panoramia}, strengthened hidden-state auditing \citep{cebere2025tighter}, studied sequential auditing \citep{gonzalez2025sequentially}, and optimized canaries for one-run auditing \citep{wang2026optimizing}. Our study is orthogonal to these directions: we keep the standard one-run setup fixed and ask whether weak-view test-time scoring helps under honest query accounting.

\section{Method}
\label{sec:method}

\subsection{Threat Model and Goal}
We assume the one-run neighboring-dataset auditing setup of \citet{steinke2023privacy}: the auditor observes only the final DP model, has no gradients or hidden states, and evaluates a fixed audit pool of member/non-member candidate pairs. All auxiliary fitting, screening, and score selection uses only non-private data. The goal is not to maximize generic attack AUC, but to maximize empirical privacy evidence, measured as an empirical \(\eps\) lower bound at \(\del=10^{-5}\).

\subsection{Public Data Splits}
For each dataset and seed, the training set is partitioned into a private base split, a public calibration split, a public proxy split, a public screening split, and an audit reserve. The executed configuration uses 42{,}000 private training examples, 6{,}000 calibration examples, 4{,}000 proxy examples, 4{,}000 screening examples, and 4{,}000 reserve examples for Fashion-MNIST; CIFAR-10 uses 32{,}000, 6{,}000, 4{,}000, 4{,}000, and 4{,}000 respectively. The audit reserve supplies a fixed pool of 240 neighboring candidate pairs per checkpoint.

\subsection{Audit Pool Construction}
Each audit pool contains 120 random in-distribution pairs and 120 FGSM-derived pairs. For every pair, one source example and one same-class reserve example define the neighboring choice. A binary membership mask determines which side is inserted into the private training set. FGSM canaries are generated with a non-private surrogate trained only on public data.

\subsection{Base Statistics and Quantile Residuals}
For an input-label pair \((x,y)\), the base statistic is the negative cross-entropy loss,
\[
b(x,y) = \log p_\theta(y \mid x),
\]
where larger values indicate higher confidence in the true label. Following \citet{liu2025enhancing}, we residualize each statistic against public difficulty features
\[
z(x) = \big(\text{surrogate loss},\ \text{surrogate entropy},\ \text{surrogate max prob},\ y,\ \text{reference feature norm}\big).
\]

For a target statistic \(t(x)\), three gradient-boosted quantile regressors estimate the 0.5, 0.9, and 0.95 conditional quantiles on the public calibration split. The residual used for auditing is
\[
r_t(x) = \frac{t(x) - \widehat{q}_{0.9}(z(x))}{\max\{\widehat{q}_{0.95}(z(x)) - \widehat{q}_{0.5}(z(x)), \tau_t\}},
\]
where \(\tau_t\) is a denominator floor learned from the public calibration distribution.

\subsection{Weak-View Statistics}
Each candidate is queried once on the clean input and \(K\) times on weak label-preserving views \(a_1(x),\dots,a_K(x)\). The executed study uses \(K=4\) for the main method and \(K=2\) for one ablation. From the weak views we compute:
\begin{align*}
\rclean(x) &= r_{b,\mathrm{clean}}(x), \\
\rmean(x) &= r_{\text{mean weak loss}}(x), \\
\rworst(x) &= r_{\text{worst weak loss}}(x), \\
\ragree(x) &= r_{\text{agreement rate}}(x).
\end{align*}

The main structured score is
\[
S_{\mathrm{structured}}(x) = \max\{\rclean(x), \rworst(x), \lambda \ragree(x)\}.
\]
We compare it against:
\begin{itemize}[leftmargin=*,topsep=2pt,itemsep=2pt]
    \item \textbf{Raw loss}: the clean negative-loss score \(b(x,y)\),
    \item \textbf{Clean quantile}: \(S_{\mathrm{clean}}(x)=\rclean(x)\),
    \item \textbf{Mean only}: \(S_{\mathrm{mean}}(x)=\rmean(x)\).
\end{itemize}

\subsection{Transform Screening}
Weak transforms are screened once per dataset on the public screening split using a public reference model. A transform passes only if clean accuracy drop is at most 1 percentage point, FGSM accuracy drop is at most 1 percentage point, clean agreement with the untransformed reference prediction is at least 98\%, and FGSM agreement is at least 98\%.

This screening is intentionally conservative and materially affects the experiment. For Fashion-MNIST, exactly two Gaussian-noise transforms survive: `noise\_001' with clean agreement 0.997 and FGSM agreement 1.000, and `noise\_003' with clean agreement 0.994 and FGSM agreement 0.992. For CIFAR-10, only one transform survives: `noise\_0005', with clean agreement 0.993 and FGSM agreement 0.988. The CIFAR-10 condition is therefore a stress case in which the \(K=4\) weak-view score effectively repeats one surviving weak transform.

\subsection{Public-Only \lambdasection{} Selection}
The scalar \(\lambda \in \{0, 0.25, 0.5, 0.75, 1.0\}\) is selected using only the public proxy split. For each dataset/seed combination, three non-private proxy auditors are trained on disjoint public pseudo-member and pseudo-nonmember subsets, then evaluated with the same matched-budget pipeline at a proxy query budget of 200. We record both the median-AUC and median-\(\eps\)-lower-bound selections, but use the \(\eps\)-selected value in the headline method.

\subsection{Matched-Budget Evaluation}
The core protocol measures privacy evidence as a function of consumed queries. If a method with query cost \(c\) per candidate audits \(m\) candidates, it consumes \(cm\) queries. For the single-view methods \(c=1\); for the \(K=4\) weak-view methods \(c=5\); for the \(K=2\) ablation \(c=3\).

For each checkpoint, we evaluate 20 random orderings of the same 240-pair audit pool over query budgets from 40 to 1200 in steps of 40. We report:
\begin{itemize}[leftmargin=*,topsep=2pt,itemsep=2pt]
    \item \textbf{Matched-budget curves}: methods are compared at equal consumed query budgets.
    \item \textbf{Same-candidate curves}: all methods are evaluated on the same prefix of candidates, allowing multi-view methods to spend more queries per candidate.
\end{itemize}

The primary metric is the empirical \(\eps\) lower bound at \(\del=10^{-5}\), computed from ROC thresholds with beta-binomial confidence bounds. Secondary metrics are ROC AUC and TPR at low FPR.

\section{Experiments}
\label{sec:experiments}

\subsection{Setup}
We evaluate on Fashion-MNIST and CIFAR-10. Fashion-MNIST uses a small CNN trained with DP-SGD for 10 epochs, batch size 512, learning rate 0.55, momentum 0.9, weight decay \(10^{-4}\), and max gradient norm 1.0. CIFAR-10 uses a frozen pretrained ResNet-18 backbone with a trained linear head, optimized with DP-SGD for 8 epochs, batch size 1024, learning rate 0.35, momentum 0.9, weight decay \(5\times 10^{-4}\), and max gradient norm 1.0. All runs use four seeds \(\{7,17,27,37\}\).

The realized privacy budgets are close to their targets: Fashion-MNIST reaches \(\eps=3.943\) and \(\eps=7.954\), while CIFAR-10 reaches \(\eps=3.942\) and \(\eps=7.914\). Mean test accuracies are 0.839 and 0.843 on Fashion-MNIST, and 0.399 and 0.402 on CIFAR-10, for targets \(\eps=4\) and \(\eps=8\) respectively. Mean per-checkpoint runtime is 1.14--1.17 minutes on Fashion-MNIST and 0.88--0.89 minutes on CIFAR-10.

\subsection{Main Results}
Table~\ref{tab:main} reports full-pool results, and Figures~\ref{fig:eps-curves} and~\ref{fig:auc-curves} show matched-budget curves. The central result is negative: the structured weak-view score never improves the empirical \(\eps\) lower bound over the clean quantile baseline under matched total-query accounting.

On CIFAR-10, the clean quantile baseline is the strongest method on the primary metric among the compared methods in the high-privacy regime and remains stronger than the structured weak-view score in both privacy settings. At target \(\eps=4\), clean quantile reaches \(0.136 \pm 0.178\) versus \(0.093 \pm 0.187\) for the structured score; at target \(\eps=8\), the gap is \(0.158 \pm 0.273\) versus \(0.078 \pm 0.155\). The mean-only weak-view control slightly exceeds clean quantile at target \(\eps=4\) on full-pool \(\eps\) lower bound (\(0.173 \pm 0.261\) versus \(0.136 \pm 0.178\)), but this advantage disappears in the matched-budget curves because it spends five queries per candidate.

On Fashion-MNIST, weak views are also unhelpful on the primary metric. At target \(\eps=4\), both clean quantile and structured score yield mean empirical lower bound \(0.000\), while raw loss has a nonzero mean of \(0.169 \pm 0.338\) driven by a single FGSM-heavy seed. At target \(\eps=8\), clean quantile attains \(0.049 \pm 0.098\), mean-only attains \(0.009 \pm 0.018\), and structured again collapses to \(0.000\). Across all four settings, the full-pool paired structured-minus-clean difference recorded in `results.json' is non-positive or zero, and no sustained positive matched-budget region appears. The secondary same-candidate curves show one small counterexample: on Fashion-MNIST at target \(\eps=4\) and budget 200, the mean structured-minus-clean difference is \(+0.0042\) (0.0141 versus 0.0099), but it disappears by budget 240 and does not carry over to the matched-budget evaluation.

\begin{table*}[t]
\caption{Full-pool audit results over all 240 audit pairs per checkpoint. All entries are mean \(\pm\) standard deviation across 4 seeds. Best \(\eps\) lower bound and AUC in each dataset/privacy block are bolded. Higher is better for all reported metrics.}
\label{tab:main}
\centering
\small
\resizebox{\textwidth}{!}{
\begin{tabular}{llcccc}
\toprule
Dataset & Target \(\eps\) & Method & Empirical \(\eps\) LB \(\uparrow\) & ROC AUC \(\uparrow\) & TPR@1\% FPR \(\uparrow\) \\
\midrule
\multirow{4}{*}{Fashion-MNIST} 
& \multirow{4}{*}{4} 
& Raw loss & \textbf{0.169 \(\pm\) 0.338} & \textbf{0.508 \(\pm\) 0.009} & 0.000 \(\pm\) 0.000 \\
& & Clean quantile & 0.000 \(\pm\) 0.000 & 0.505 \(\pm\) 0.014 & \textbf{0.023 \(\pm\) 0.018} \\
& & Mean only & 0.000 \(\pm\) 0.000 & 0.506 \(\pm\) 0.010 & 0.019 \(\pm\) 0.014 \\
& & Structured weak views & 0.000 \(\pm\) 0.000 & 0.496 \(\pm\) 0.023 & 0.014 \(\pm\) 0.019 \\
\midrule
\multirow{4}{*}{Fashion-MNIST} 
& \multirow{4}{*}{8} 
& Raw loss & \textbf{0.052 \(\pm\) 0.103} & 0.459 \(\pm\) 0.019 & 0.000 \(\pm\) 0.000 \\
& & Clean quantile & 0.049 \(\pm\) 0.098 & 0.463 \(\pm\) 0.022 & 0.016 \(\pm\) 0.017 \\
& & Mean only & 0.009 \(\pm\) 0.018 & 0.467 \(\pm\) 0.020 & \textbf{0.022 \(\pm\) 0.018} \\
& & Structured weak views & 0.000 \(\pm\) 0.000 & \textbf{0.497 \(\pm\) 0.008} & 0.018 \(\pm\) 0.014 \\
\midrule
\multirow{4}{*}{CIFAR-10} 
& \multirow{4}{*}{4} 
& Raw loss & 0.011 \(\pm\) 0.022 & 0.484 \(\pm\) 0.029 & 0.012 \(\pm\) 0.005 \\
& & Clean quantile & 0.136 \(\pm\) 0.178 & 0.496 \(\pm\) 0.036 & 0.027 \(\pm\) 0.027 \\
& & Mean only & \textbf{0.173 \(\pm\) 0.261} & 0.496 \(\pm\) 0.038 & \textbf{0.031 \(\pm\) 0.025} \\
& & Structured weak views & 0.093 \(\pm\) 0.187 & \textbf{0.503 \(\pm\) 0.012} & 0.029 \(\pm\) 0.030 \\
\midrule
\multirow{4}{*}{CIFAR-10} 
& \multirow{4}{*}{8} 
& Raw loss & 0.008 \(\pm\) 0.012 & 0.476 \(\pm\) 0.027 & 0.012 \(\pm\) 0.005 \\
& & Clean quantile & \textbf{0.158 \(\pm\) 0.273} & 0.481 \(\pm\) 0.030 & 0.012 \(\pm\) 0.008 \\
& & Mean only & 0.146 \(\pm\) 0.280 & 0.481 \(\pm\) 0.026 & \textbf{0.014 \(\pm\) 0.010} \\
& & Structured weak views & 0.078 \(\pm\) 0.155 & \textbf{0.499 \(\pm\) 0.022} & \textbf{0.014 \(\pm\) 0.010} \\
\bottomrule
\end{tabular}
}
\end{table*}

\begin{figure*}[t]
\centering
\includegraphics[width=0.9\linewidth]{figures/figure1_matched_budget_eps.pdf}
\caption{Matched-budget empirical \(\eps\) lower-bound curves versus consumed queries. Means and standard deviations are aggregated over 4 seeds and 20 random audit-pool orderings per seed. Across all four dataset/privacy settings, the structured weak-view curve does not show a sustained advantage over the clean quantile baseline.}
\label{fig:eps-curves}
\end{figure*}

\begin{figure*}[t]
\centering
\includegraphics[width=0.9\linewidth]{figures/figure2_matched_budget_auc.pdf}
\caption{Matched-budget ROC AUC curves versus consumed queries. Weak-view aggregation can mildly improve AUC, especially on CIFAR-10, but these AUC gains do not translate into stronger empirical \(\eps\) lower bounds once the extra query cost is accounted for.}
\label{fig:auc-curves}
\end{figure*}

\subsection{Ablations}
Figure~\ref{fig:ablations} studies which parts of the weak-view pipeline matter. The ablations confirm the main result rather than overturning it.

First, residualization matters. On CIFAR-10 at target \(\eps=4\), removing residualization drops the mean full-pool empirical lower bound of the structured score from 0.093 to 0.000. On Fashion-MNIST at target \(\eps=8\), removing residualization raises the mean lower bound from 0.000 to 0.047, but still does not beat the clean quantile baseline at 0.049.

Second, the agreement feature is mostly inactive because public proxy selection chooses \(\lambda_{\eps}=0\) in every seed and every setting. This means the headline structured score reduces to \(\max\{\rclean,\rworst\}\) in the executed experiments. The \(\lambda\)-selection pathology is visible in Figure~\ref{fig:lambda}: the median proxy empirical lower bound is 0.0 for every \(\lambda\), while proxy AUC weakly favors \(\lambda=0.25\).

Third, the \(\lambda\)-selection criterion is itself unstable. The \(\eps\)-tuned and AUC-tuned selections agree in all four Fashion-MNIST target-\(\eps=4\) seeds, disagree in one of four Fashion-MNIST target-\(\eps=8\) seeds, and disagree in all eight CIFAR-10 seeds across the two privacy regimes. In those CIFAR-10 runs, \(\lambda_{\eps}=0\) for every seed while \(\lambda_{\mathrm{AUC}}=0.25\) for every seed. This matters for robustness: the agreement feature is not supported by the privacy objective even when it modestly improves a proxy attack metric, so the method's composition is sensitive to the tuning target rather than stably identified by the public proxy audits.

Fourth, reducing the number of weak views does not rescue the method. The \(K=2\) structured variant remains below or equal to the clean quantile baseline on the primary metric in all four settings.

\begin{figure*}[t]
\centering
\includegraphics[width=0.9\linewidth]{figures/figure3_ablations.pdf}
\caption{Ablation summary. Left: full-method and component-removal comparisons across dataset/privacy settings. Right: breakdown by canary family. The structured weak-view score does not deliver stronger empirical \(\eps\) lower bounds than the clean quantile baseline, and the CIFAR-10 stress case highlights that only one screened transform survives.}
\label{fig:ablations}
\end{figure*}

\begin{figure}[t]
\centering
\includegraphics[width=0.8\linewidth]{figures/figure4_lambda_selection.pdf}
\caption{Proxy \(\lambda\)-selection summary aggregated across all settings. Median proxy empirical \(\eps\) lower bound is 0.0 for every \(\lambda\), while proxy AUC is slightly higher for \(\lambda \ge 0.25\) than for \(\lambda=0\). In the executed experiments, \(\lambda_{\eps}\) selects 0.0 in all 16 dataset/seed/privacy combinations.}
\label{fig:lambda}
\end{figure}

\subsection{Analysis}
Two observations explain the negative result. First, the matched-budget accounting is strict: a \(K=4\) weak-view score spends five queries per candidate, so any per-candidate gain must be large to beat a calibrated single-view method at equal query cost. The same-candidate curves are mostly negative as well; Appendix~\ref{app:same-candidate} shows only one positive mean structured-minus-clean difference, \(+0.0042\) on Fashion-MNIST at target \(\eps=4\) and budget 200, and even that exception vanishes by budget 240.

Second, the weak transforms that survive screening are extremely conservative. Fashion-MNIST retains only Gaussian noise with standard deviations 0.01 and 0.03, and CIFAR-10 retains only Gaussian noise with standard deviation 0.005. The local neighborhoods explored by the weak views are therefore narrow, especially on CIFAR-10, where repeated views mostly probe one near-identity transform. Under those conditions, extra queries appear largely redundant once clean-view difficulty has already been calibrated.

\section{Discussion and Limitations}
\label{sec:discussion}
The main implication is practical rather than algorithmic. In this one-run final-model-only setting, researchers should not assume that repeated weak views automatically strengthen empirical privacy auditing. If a strong single-view quantile baseline is already available, extra views may increase AUC slightly yet still fail to improve the privacy metric that matters after honest accounting.

The study has several limitations. First, it is intentionally compute bounded: we evaluate only two datasets, two privacy targets, and four seeds. Second, the CIFAR-10 model is a linear probe on a frozen pretrained ResNet-18, not full end-to-end DP training. Third, transform screening is conservative by design, and on CIFAR-10 only one transform survives, making the weak-view condition unusually narrow. Fourth, proxy \(\lambda\)-selection is degenerate in this regime: the mean fraction of zero proxy empirical lower bounds is 0.939 on Fashion-MNIST at target \(\eps=4\), 0.954 on Fashion-MNIST at target \(\eps=8\), 0.951 on CIFAR-10 at target \(\eps=4\), and 0.958 on CIFAR-10 at target \(\eps=8\). This weakens our ability to test richer agreement-weighted combinations.

These limitations also point to future work. More diverse but still label-preserving transforms, stronger proxy selection objectives, and settings with larger audit pools or stronger base leakage may reveal regimes where weak views help. Our result is therefore not ``weak views never help,'' but ``they do not help here once the comparison is made against a strong calibrated baseline under matched total-query accounting.''

\section{Conclusion}
\label{sec:conclusion}
We evaluated whether weak-view residual scores improve one-run DP auditing beyond a quantile-calibrated single-view baseline. Using fixed audit pools, public-only calibration and tuning, and matched-total-query evaluation on Fashion-MNIST and CIFAR-10, we found that they do not. The structured weak-view score never exceeded the clean quantile baseline on the primary metric, the empirical \(\eps\) lower bound at \(\del=10^{-5}\), in any of the four dataset/privacy settings; on CIFAR-10, for example, clean quantile achieved mean lower bounds of 0.136 and 0.158 at target \(\eps=4\) and \(\eps=8\), compared with 0.093 and 0.078 for the structured score.

This is a negative result, but a useful one. It suggests that in compute-bounded one-run auditing, effort may be better spent on stronger single-view calibration and better audit-pool construction than on additional weak-view queries that consume budget without improving privacy evidence. Future work should test whether this conclusion changes under richer transform families, larger audits, or stronger proxy selection signals.

\appendix

\section{Reproducibility Appendix}
\label{app:repro}

The executed environment differed from the preregistered stack. The run used the preexisting Conda environment `/home/zz865/.conda/envs/ar' rather than the planned Python 3.10 stack; the saved `version\_note' records Python 3.11, Opacus 1.5.4, NumPy 2.x, and a CIFAR-10 preprocessing correction to ImageNet normalization for the frozen pretrained ResNet-18 backbone. The saved hardware snapshot reports one NVIDIA RTX A6000 GPU (driver 570.124.06, CUDA 12.8), 4 CPU cores, and 503~GiB RAM. The full package list is stored in `results.json' and `exp/environment/logs/system\_info.json'.

All stochastic components used the executed seeds \(\{7,17,27,37\}\), with audit-ordering seeds \(\{0,\dots,19\}\), query budgets \(\{40,80,\dots,1200\}\), audit-pool size 240, primary weak-view count \(K=4\), and \(K=2\) as the only executed weak-view ablation. The reporting metadata explicitly records two protocol deviations: the environment mismatch above and the proposal-to-execution seed-count mismatch, since `proposal.md' still mentions 3 seeds while the executed study uses 4. Two optional plan items were not run: the true-label margin sensitivity analysis and the \(K=8\) weak-view ablation.

Transform screening used canonical seed 7 and produced the following frozen recipes. On Fashion-MNIST, only `noise\_001' and `noise\_003' passed screening, with \((\mathrm{clean\ agree}, \mathrm{FGSM\ agree})=(0.997,1.000)\) and \((0.994,0.992)\) respectively; the dataset is not marked as a stress case. On CIFAR-10, only `noise\_0005' passed, with \((0.993,0.988)\), and the saved screening log marks the dataset as a stress case.

\begin{table}[t]
\caption{Executed \(\lambda\)-selection outcomes by setting. Entries list the per-seed public-proxy selections for the privacy objective and the AUC objective, in seed order \(\{7,17,27,37\}\).}
\label{tab:lambda-appendix}
\centering
\small
\begin{tabular}{lll}
\toprule
Setting & \(\lambda_{\eps}\) & \(\lambda_{\mathrm{AUC}}\) \\
\midrule
Fashion-MNIST, target \(\eps=4\) & [0, 0, 0, 0] & [0, 0, 0, 0] \\
Fashion-MNIST, target \(\eps=8\) & [0, 0, 0, 0] & [0, 0, 0.25, 0] \\
CIFAR-10, target \(\eps=4\) & [0, 0, 0, 0] & [0.25, 0.25, 0.25, 0.25] \\
CIFAR-10, target \(\eps=8\) & [0, 0, 0, 0] & [0.25, 0.25, 0.25, 0.25] \\
\bottomrule
\end{tabular}
\end{table}

\section{Same-Candidate Summary}
\label{app:same-candidate}

Table~\ref{tab:same-candidate-diff} summarizes the secondary same-candidate analysis by reporting the mean empirical \(\eps\) lower-bound difference \(S_{\mathrm{structured}}-S_{\mathrm{clean}}\) at representative budgets. The one positive entry is Fashion-MNIST at target \(\eps=4\) and budget 200; all other reported entries are non-positive, and the difference is 0.0000 or negative by budget 240 in every setting.

\begin{table}[t]
\caption{Secondary same-candidate mean empirical \(\eps\) lower-bound differences, computed as structured minus clean quantile from `artifacts/figures/source\_data/same\_candidate\_curves.csv'. Positive values favor the structured weak-view score.}
\label{tab:same-candidate-diff}
\centering
\small
\begin{tabular}{lcccc}
\toprule
Setting & 120 & 160 & 200 & 240 \\
\midrule
Fashion-MNIST, target \(\eps=4\) & -0.0012 & -0.0076 & 0.0042 & 0.0000 \\
Fashion-MNIST, target \(\eps=8\) & -0.0542 & -0.0726 & -0.0638 & -0.0488 \\
CIFAR-10, target \(\eps=4\) & -0.0297 & -0.0331 & -0.0318 & -0.0430 \\
CIFAR-10, target \(\eps=8\) & -0.0475 & -0.0412 & -0.0721 & -0.0809 \\
\bottomrule
\end{tabular}
\end{table}

\bibliographystyle{plainnat}
\bibliography{references}

\end{document}
